\chapter{《Attention Is All You Need》中文版}

\section{摘要}

目前主流的序列转换模型建立在复杂的循环神经网络或卷积神经网络之上，通常包含一个编码器和一个解码器。性能最优的模型还通过注意力机制将编码器与解码器连接起来。本文提出一种全新的简单网络架构——Transformer，它完全基于注意力机制，彻底摒弃了循环结构与卷积操作。在两项机器翻译任务上的实验表明，该模型在翻译质量上更优，同时具备更强的并行性，所需训练时间也大幅缩短。在 WMT 2014 英德翻译任务上，我们的模型取得了 28.4 的 BLEU 分数，比现有最优结果（包括集成模型）提升了 2 个 BLEU 以上。在 WMT 2014 英法翻译任务上，我们的模型仅用 8 块 GPU 训练 3.5 天，便以单模型创下 41.8 的 BLEU 新纪录，训练成本仅为文献中最优模型的一小部分。此外，我们将 Transformer 成功应用于英语成分句法分析任务（无论训练数据规模大小），证明其具备良好的泛化能力。

\section{引言}

循环神经网络，尤其是长短期记忆网络（LSTM）\cite{hochreiter1997long}和门控循环神经网络\cite{cho2014learning}，已在序列建模与转换问题（如语言建模和机器翻译）中牢固确立了其最先进方法的地位\cite{sutskever2014sequence,cho2014learning,bahdanau2014neural}。此后，大量研究持续推动着循环语言模型与编码器-解码器架构的边界\cite{wu2016google,luong2015effective,jozefowicz2016exploring}。

循环模型通常沿输入和输出序列的符号位置逐步推进计算。将位置与计算时间步对齐后，模型以上一时刻的隐藏状态 $h_{t-1}$ 和当前位置 $t$ 的输入为函数，生成一系列隐藏状态 $h_t$。这种本质上的顺序性使得训练样本内部无法并行化，在序列较长时尤为制约性能，因为内存限制会压缩跨样本的批处理规模。近期工作通过分解技巧\cite{kuchaiev2017factorization}和条件计算\cite{shazeer2017outrageously}在计算效率上取得了显著提升，后者同时也改善了模型性能。然而，顺序计算这一根本约束依然存在。

注意力机制已成为各类序列建模与转换模型中不可或缺的组成部分，它允许建模输入或输出序列中任意距离的依赖关系\cite{bahdanau2014neural,structuredattentionnetworks}。然而，除少数例外\cite{decomposableattentionmodel}，注意力机制几乎都与循环网络配合使用。

本文提出 Transformer——一种完全摒弃循环结构、仅依赖注意力机制来捕捉输入与输出之间全局依赖关系的模型架构。Transformer 具备更强的并行化能力，仅需在 8 块 P100 GPU 上训练短短十二小时，便可在翻译质量上达到新的最先进水平。

\section{背景}

减少顺序计算这一目标，同样是 Extended Neural GPU\cite{extendedneuralgpu}、ByteNet\cite{kalchbrenner2016neural} 和 ConvS2S\cite{gehring2017convolutional} 的出发点。这些模型均以卷积神经网络为基本构件，能够并行计算所有输入和输出位置的隐藏表示。但在这些模型中，关联两个任意输入或输出位置信号所需的操作次数，随位置间距的增大而增长——ConvS2S 呈线性增长，ByteNet 呈对数增长。这使得学习远距离位置之间的依赖关系更加困难\cite{hochreiter2001gradient}。在 Transformer 中，这一代价被降低为常数次操作，尽管代价是因对注意力权重位置取平均而导致有效分辨率下降——我们通过第 3.2 节介绍的多头注意力机制来抵消这一影响。

自注意力（有时称为内部注意力）是一种将单个序列中不同位置相互关联、从而计算序列表示的注意力机制。自注意力已在阅读理解、抽象摘要、文本蕴含及与任务无关的句子表示学习等多种任务中取得成功\cite{cheng2016long,decomposableattentionmodel,paulus2017deep,lin2017structured}。

端到端记忆网络基于循环注意力机制（而非序列对齐的循环结构），已被证明在简单语言问答和语言建模任务上表现良好\cite{sukhbaatar2015}。

然而，据我们所知，Transformer 是首个完全依赖自注意力来计算输入和输出表示、不使用任何序列对齐 RNN 或卷积的转换模型。在接下来的章节中，我们将描述 Transformer，阐述自注意力的动机，并讨论其相较于\cite{sutskever2014sequence,bahdanau2014neural}等模型的优势。

\section{模型架构}

大多数具有竞争力的神经序列转换模型都采用编码器-解码器结构\cite{cho2014learning,bahdanau2014neural,sutskever2014sequence}。编码器将输入符号表示序列 $(x_1, \ldots, x_n)$ 映射为连续表示序列 $\mathbf{z} = (z_1, \ldots, z_n)$。给定 $\mathbf{z}$ 后，解码器逐步生成输出符号序列 $(y_1, \ldots, y_m)$，每次生成一个元素。在每一步中，模型是自回归的\cite{graves2013generating}，即在生成下一个符号时，将已生成的符号作为额外输入。

Transformer 遵循上述整体架构，对编码器和解码器均采用堆叠的自注意力层与逐点全连接层，分别对应图 1 的左半部分和右半部分。

\subsection{编码器与解码器堆栈}

\textbf{编码器：}编码器由 $N = 6$ 个相同层堆叠而成。每层包含两个子层：第一个是多头自注意力机制，第二个是简单的逐位置全连接前馈网络。在每个子层周围采用残差连接\cite{he2016deep}，随后进行层归一化\cite{ba2016layer}。即每个子层的输出为 $\mathrm{LayerNorm}(x + \mathrm{Sublayer}(x))$，其中 $\mathrm{Sublayer}(x)$ 为该子层自身实现的函数。为便于使用残差连接，模型中所有子层及嵌入层的输出维度均为 $d_{\text{model}} = 512$。

\textbf{解码器：}解码器同样由 $N = 6$ 个相同层堆叠而成。在编码器每层的两个子层基础上，解码器额外插入第三个子层，对编码器堆栈的输出执行多头注意力。与编码器类似，每个子层周围同样采用残差连接并进行层归一化。此外，对解码器堆栈中的自注意力子层进行了修改，通过掩码（Masking）防止当前位置关注到后续位置。这种掩码结合输出嵌入偏移一个位置的设计，确保位置 $i$ 的预测只能依赖位置 $i$ 之前的已知输出。

\subsection{注意力}

注意力函数可描述为：将一个查询（query）和一组键值对（key-value pairs）映射为输出，其中查询、键、值和输出均为向量。输出是值的加权求和，每个值的权重由查询与对应键的相容性函数计算得到。

\subsubsection{缩放点积注意力}

我们将本文的注意力称为"缩放点积注意力"（图 2）。输入包含维度为 $d_k$ 的查询和键，以及维度为 $d_v$ 的值。将查询与所有键做点积，除以 $\sqrt{d_k}$ 后，应用 softmax 函数得到各值的权重。

在实践中，我们将一组查询打包成矩阵 $Q$，同时将键和值分别打包成矩阵 $K$ 和 $V$，并行计算注意力函数，输出矩阵为：
\[
  \mathrm{Attention}(Q, K, V) = \mathrm{softmax}\!\left(\frac{QK^{\top}}{\sqrt{d_k}}\right)V
\]

最常用的两种注意力函数是加性注意力\cite{bahdanau2014neural}和点积（乘法）注意力。点积注意力与本文算法完全相同，区别仅在于缩放因子 $\frac{1}{\sqrt{d_k}}$。加性注意力使用单隐层前馈网络计算相容性函数。二者理论复杂度相近，但点积注意力在实践中速度更快、空间效率更高，因为可以借助高度优化的矩阵乘法代码实现。

当 $d_k$ 较小时，两种机制性能相当；当 $d_k$ 较大时，不加缩放的点积注意力性能逊于加性注意力\cite{DBLP:journals/corr/BritzGLL17}。我们推测，当 $d_k$ 较大时，点积结果幅值增大，将 softmax 函数推入梯度极小的区域。为抵消这一效应，我们将点积结果除以 $\frac{1}{\sqrt{d_k}}$ 进行缩放。

\subsubsection{多头注意力}

我们发现，与其用 $d_{\text{model}}$ 维的键、值和查询执行单次注意力函数，不如将查询、键和值分别经过 $h$ 次不同的可学习线性投影，投影到 $d_k$、$d_k$ 和 $d_v$ 维度，效果更佳。在每组投影后的查询、键和值上并行执行注意力函数，得到 $d_v$ 维输出值，将其拼接后再次投影，得到最终输出，如图 2 所示。

多头注意力使模型能够在不同位置上同时关注来自不同表示子空间的信息，而单头注意力的平均化操作会抑制这一能力。
\[
  \mathrm{MultiHead}(Q, K, V) = \mathrm{Concat}(\mathrm{head}_1, \ldots, \mathrm{head}_h)\,W^O
\]
其中 $\mathrm{head}_i = \mathrm{Attention}(QW_i^Q,\, KW_i^K,\, VW_i^V)$

投影矩阵参数为
$W_i^Q \in \mathbb{R}^{d_{\text{model}} \times d_k}$，
$W_i^K \in \mathbb{R}^{d_{\text{model}} \times d_k}$，
$W_i^V \in \mathbb{R}^{d_{\text{model}} \times d_v}$，
$W^O \in \mathbb{R}^{hd_v \times d_{\text{model}}}$。

本文采用 $h = 8$ 个并行注意力头，每个头的维度为 $d_k = d_v = d_{\text{model}}/h = 64$。由于每个头的维度缩减，总计算量与全维度单头注意力相近。

\subsubsection{注意力在模型中的应用}

Transformer 以三种不同方式使用多头注意力：

\begin{itemize}
  \item 在"编码器-解码器注意力"层中，查询来自上一解码器层，键和值来自编码器的输出。这使解码器的每个位置都能关注输入序列的所有位置，与 seq2seq 模型中的典型编码器-解码器注意力机制\cite{wu2016google,bahdanau2014neural,decomposableattentionmodel}类似。
  \item 编码器包含自注意力层。在自注意力层中，键、值和查询均来自同一处——上一编码器层的输出。编码器中每个位置都可以关注编码器上一层的所有位置。
  \item 同样，解码器中的自注意力层允许每个位置关注解码器中该位置及之前的所有位置。为保持自回归性质，需要阻止解码器中信息向左流动。我们在缩放点积注意力内部实现这一点：将 softmax 输入中对应非法连接的所有值掩码为 $-\infty$。
\end{itemize}

\subsection{逐位置前馈网络}

除注意力子层外，编码器和解码器的每一层还包含一个全连接前馈网络，对每个位置分别且相同地应用。该网络由两个线性变换组成，中间加入 ReLU 激活函数：
\[
  \mathrm{FFN}(x) = \max(0,\, xW_1 + b_1)\,W_2 + b_2
\]
不同位置的线性变换相同，但不同层之间使用不同参数。也可将其理解为两个核大小为 1 的卷积。输入和输出维度为 $d_{\text{model}} = 512$，内层维度为 $d_{ff} = 2048$。

\subsection{嵌入与 Softmax}

与其他序列转换模型类似，我们使用可学习的嵌入将输入词元和输出词元转换为 $d_{\text{model}}$ 维向量。同时使用常规的可学习线性变换和 softmax 函数，将解码器输出转换为预测下一词元的概率。本模型在两个嵌入层与 pre-softmax 线性变换之间共享权重矩阵\cite{press2016using}，在嵌入层中将权重乘以 $\sqrt{d_{\text{model}}}$。

\subsection{位置编码}

由于本模型不含循环结构也不含卷积，为使模型能够利用序列的顺序信息，必须向模型注入词元在序列中的相对或绝对位置信息。为此，我们在编码器和解码器堆栈底部的输入嵌入上叠加"位置编码"。位置编码与嵌入的维度相同，均为 $d_{\text{model}}$，因此二者可以直接相加。位置编码有多种选择，包括可学习的和固定的\cite{gehring2017convolutional}。

本文使用不同频率的正弦和余弦函数：
\[
  PE_{(pos,\,2i)}   = \sin\!\left(\frac{pos}{10000^{2i/d_{\text{model}}}}\right)
  \qquad
  PE_{(pos,\,2i+1)} = \cos\!\left(\frac{pos}{10000^{2i/d_{\text{model}}}}\right)
\]
其中 $pos$ 为位置，$i$ 为维度索引。每个维度的位置编码对应一条正弦曲线，波长从 $2\pi$ 到 $10000 \cdot 2\pi$ 构成等比数列。选择该函数是因为我们假设它能使模型轻松学会通过相对位置进行注意力计算——对于任意固定偏移 $k$，$PE_{pos+k}$ 可表示为 $PE_{pos}$ 的线性函数。

我们还实验了使用可学习的位置嵌入\cite{gehring2017convolutional}，发现两种方式结果几乎相同（见表 3 第 (E) 行）。最终选择正弦版本，因为它可能允许模型外推至训练时未见过的更长序列。

\section{为什么使用自注意力}

本节将自注意力层与循环层、卷积层在多个方面进行比较。这些层通常用于将一段变长符号表示序列 $(x_1, \ldots, x_n)$ 映射为另一段等长序列 $(z_1, \ldots, z_n)$，其中 $x_i, z_i \in \mathbb{R}^d$，例如典型序列转换编码器或解码器中的隐藏层。我们从三个方面论证使用自注意力的动机。

其一是每层的总计算复杂度。其二是可并行化的计算量，以所需最少顺序操作数衡量。

其三是网络中长距离依赖之间的路径长度。学习长距离依赖是许多序列转换任务的核心挑战。影响这一能力的关键因素之一，是前向与反向信号在网络中需要经过的路径长度。输入与输出序列中任意位置组合之间的路径越短，越容易学习长距离依赖\cite{hochreiter2001gradient}。因此，我们也比较了由不同类型层构成的网络中，任意两个输入输出位置之间的最大路径长度。

如表 1 所示，自注意力层以常数次顺序操作连接所有位置，而循环层需要 $O(n)$ 次顺序操作。在计算复杂度方面，当序列长度 $n$ 小于表示维度 $d$ 时，自注意力层比循环层更快——这在机器翻译中最先进模型所使用的句子表示（如词片\cite{wu2016google}和字节对\cite{sennrich2015neural}表示）中最为常见。

为提升处理超长序列任务时的计算性能，自注意力可限制为只关注以各输出位置为中心、大小为 $r$ 的邻域。这将使最大路径长度增至 $O(n/r)$，我们计划在未来工作中进一步研究这一方法。

核宽为 $k < n$ 的单层卷积层无法连接所有输入输出位置对。若要做到这一点，连续核情况下需要堆叠 $O(n/k)$ 层卷积，膨胀卷积\cite{yu2015multi}情况下需要 $O(\log_k(n))$ 层，这都会增加网络中任意两个位置之间最长路径的长度。卷积层的计算代价通常比循环层高出 $k$ 倍。然而，可分离卷积\cite{chollet2016xception}可将复杂度大幅降低至 $O(k \cdot n \cdot d + n \cdot d^2)$。即便取 $k = n$，可分离卷积的复杂度也与自注意力层加逐点前馈层的组合相当——正是本文模型所采用的方案。

自注意力还有一个附带好处：可产生更具可解释性的模型。我们对模型的注意力分布进行了检视，并在附录中列举和讨论了若干示例。不仅各个注意力头明显学会了执行不同的任务，许多头还表现出与句子句法和语义结构相关的行为。

\section{训练}

本节介绍模型的训练方案。

\subsection{训练数据与批处理}

我们在标准的 WMT 2014 英德数据集上进行训练，该数据集约包含 450 万个句子对。句子采用字节对编码\cite{sennrich2015neural}，源语言与目标语言共享约 37000 个词元的词表。对于英法翻译，我们使用规模更大的 WMT 2014 英法数据集，共 3600 万个句子，词元切分为 32000 个词片词表\cite{wu2016google}。句子对按近似序列长度打包成批次，每个训练批次约包含 25000 个源词元和 25000 个目标词元。

\subsection{硬件与训练计划}

我们在一台配备 8 块 NVIDIA P100 GPU 的机器上训练模型。使用本文所描述超参数的基础模型，每个训练步耗时约 0.4 秒，共训练 100000 步，历时约 12 小时。大模型每步耗时 1.0 秒，共训练 300000 步，历时 3.5 天。

\subsection{优化器}

我们使用 Adam 优化器\cite{kingma2014adam}，参数设置为 $\beta_1 = 0.9$，$\beta_2 = 0.98$，$\epsilon = 10^{-9}$。训练过程中按以下公式动态调整学习率：
\[
  lrate = d_{\text{model}}^{-0.5} \cdot
  \min\!\left(step\_num^{-0.5},\;
  step\_num \cdot warmup\_steps^{-1.5}\right)
\]
即在前 $warmup\_steps$ 步线性增大学习率，之后按步数的平方根倒数衰减。本文设 $warmup\_steps = 4000$。

\subsection{正则化}

训练中采用三种正则化方式：

\textbf{残差 Dropout。}在每个子层的输出加到子层输入并归一化之前，对其施加 Dropout\cite{srivastava2014dropout}。此外，在编码器和解码器堆栈中，对嵌入与位置编码之和也施加 Dropout。基础模型的丢弃率为 $P_{drop} = 0.1$。

\textbf{标签平滑。}训练时采用值为 $\epsilon_{ls} = 0.1$ 的标签平滑\cite{szegedy2015rethinking}。这会使模型变得更加"不确定"，从而损害困惑度，但能提升准确率和 BLEU 分数。

\section{实验结果}

\subsection{机器翻译}

在 WMT 2014 英德翻译任务上，大型 Transformer 模型（表 2 中的 Transformer (big)）以超过 2.0 BLEU 的优势超越了此前所有已报告的最优模型（包括集成模型），以 28.4 的 BLEU 分数创下新的最先进水平。该模型的配置见表 3 最后一行，在 8 块 P100 GPU 上训练耗时 3.5 天。即便是我们的基础模型，也以远低于任何竞争模型的训练成本超越了此前所有已发表的模型和集成模型。

在 WMT 2014 英法翻译任务上，我们的大模型取得了 41.0 的 BLEU 分数，超越了此前所有已发表的单模型，训练成本不足此前最优模型的四分之一。英法任务的大型 Transformer 模型使用的 Dropout 率为 $P_{drop} = 0.1$，而非 0.3。

对于基础模型，我们对最后 5 个检查点（每 10 分钟保存一次）取平均，得到单一模型。对于大模型，则对最后 20 个检查点取平均。推理时使用束搜索，束大小为 4，长度惩罚系数 $\alpha = 0.6$\cite{wu2016google}，这些超参数均经过开发集实验选定。推理时最大输出长度设为输入长度加 50，但在可能时提前终止\cite{wu2016google}。

表 2 汇总了我们的实验结果，并将翻译质量和训练成本与文献中其他模型架构进行了比较。训练模型所用的浮点运算次数，由训练时间、所用 GPU 数量以及每块 GPU 的单精度浮点运算峰值估算得出。

\subsection{模型变体}

为评估 Transformer 各组件的重要性，我们以不同方式对基础模型进行变体实验，在开发集 newstest2013 上衡量英德翻译性能的变化，结果见表 3。推理使用上节所述的束搜索，但不进行检查点平均。

在表 3 第 (A) 组，我们按第 3.2.2 节的描述，在保持计算量不变的前提下，改变注意力头数及注意力键值维度。结果显示，单头注意力比最优设置差 0.9 BLEU，但头数过多时质量同样下降。

在第 (B) 组，我们观察到减小注意力键维度 $d_k$ 会损害模型质量，这表明判断相容性并不容易，比点积更复杂的相容性函数或许会更有益。在第 (C) 和 (D) 组，进一步验证了更大的模型效果更好，Dropout 对于防止过拟合非常有效。在第 (E) 行，用可学习的位置嵌入\cite{gehring2017convolutional}替换正弦位置编码后，结果与基础模型几乎相同。

\subsection{英语成分句法分析}

为评估 Transformer 能否泛化到其他任务，我们在英语成分句法分析任务上进行了实验。该任务有其特殊挑战：输出受到强结构约束，且通常远长于输入；此外，在小数据场景下，RNN seq2seq 模型难以达到最先进水平\cite{vinyals2015grammar}。

我们在 Penn Treebank\cite{marcus1993building} 的华尔街日报（WSJ）部分（约 4 万训练句）上训练了一个 4 层、$d_{\text{model}} = 1024$ 的 Transformer。同时还在半监督设置下进行了训练，利用了约 1700 万句的高置信度语料和 BerkeleyParser 语料\cite{vinyals2015grammar,zhu2013fast}。纯 WSJ 设置使用 16K 词元词表，半监督设置使用 32K 词元词表。

我们仅在第 22 节开发集上进行了少量实验以选择 Dropout 率（注意力和残差，见第 5.4 节）、学习率和束大小，其余所有参数均沿用英德基础翻译模型的设置。推理时最大输出长度设为输入长度加 300，束大小为 21，$\alpha = 0.3$（纯 WSJ 和半监督设置相同）。

表 4 的结果表明，尽管缺乏针对特定任务的调优，我们的模型表现出乎意料地好，优于此前所有已报告模型，仅次于循环神经网络文法模型（RNNG）\cite{dyer2016recurrent}。与 RNN seq2seq 模型\cite{vinyals2015grammar}不同，即使仅在 4 万句的 WSJ 训练集上训练，Transformer 也超越了 BerkeleyParser\cite{petrov2006learning}。

\section{结论}

本文提出了 Transformer——首个完全基于注意力机制的序列转换模型，以多头自注意力取代了编码器-解码器架构中最常用的循环层。

在翻译任务上，Transformer 的训练速度显著快于基于循环或卷积层的架构。在 WMT 2014 英德和英法翻译任务上均取得了新的最先进水平，前者甚至超越了此前所有已报告的集成模型。

我们对基于注意力模型的未来充满期待，并计划将其应用于其他任务。我们还计划将 Transformer 扩展到文本以外的输入输出模态，并研究局部受限注意力机制，以高效处理图像、音频和视频等大规模输入输出。减少生成过程的顺序性也是我们的研究目标之一。

本文训练和评估模型所用的代码已开源，地址为 \url{https://github.com/tensorflow/tensor2tensor}。

\chapter{Attention Is All You Need}

\section{Abstract}

The dominant sequence transduction models are based on complex recurrent or
convolutional neural networks that include an encoder and a decoder. The best
performing models also connect the encoder and decoder through an attention
mechanism. We propose a new simple network architecture, the Transformer,
based solely on attention mechanisms, dispensing with recurrence and
convolutions entirely. Experiments on two machine translation tasks show these
models to be superior in quality while being more parallelizable and requiring
significantly less time to train. Our model achieves 28.4 BLEU on the WMT
2014 English-to-German translation task, improving over the existing best
results, including ensembles, by over 2 BLEU. On the WMT 2014
English-to-French translation task, our model establishes a new single-model
state-of-the-art BLEU score of 41.8 after training for 3.5 days on eight
GPUs, a small fraction of the training costs of the best models from the
literature. We show that the Transformer generalizes well to other tasks by
applying it successfully to English constituency parsing both with large and
limited training data.

\section{Introduction}

Recurrent neural networks, long short-term memory \cite{hochreiter1997long}
and gated recurrent \cite{cho2014learning} neural networks in particular, have
been firmly established as state of the art approaches in sequence modeling
and transduction problems such as language modeling and machine translation
\cite{sutskever2014sequence,cho2014learning,bahdanau2014neural}. Numerous
efforts have since continued to push the boundaries of recurrent language
models and encoder-decoder architectures
\cite{wu2016google,luong2015effective,jozefowicz2016exploring}.

Recurrent models typically factor computation along the symbol positions of
the input and output sequences. Aligning the positions to steps in computation
time, they generate a sequence of hidden states $h_t$, as a function of the
previous hidden state $h_{t-1}$ and the input for position $t$. This inherently
sequential nature precludes parallelization within training examples, which
becomes critical at longer sequence lengths, as memory constraints limit
batching across examples. Recent work has achieved significant improvements in
computational efficiency through factorization tricks
\cite{kuchaiev2017factorization} and conditional computation
\cite{shazeer2017outrageously}, while also improving model performance in case
of the latter. The fundamental constraint of sequential computation, however,
remains.

Attention mechanisms have become an integral part of compelling sequence
modeling and transduction models in various tasks, allowing modeling of
dependencies without regard to their distance in the input or output sequences
\cite{bahdanau2014neural,structuredattentionnetworks}. In all but a few cases
\cite{decomposableattentionmodel}, however, such attention mechanisms are used
in conjunction with a recurrent network.

In this work we propose the Transformer, a model architecture eschewing
recurrence and instead relying entirely on an attention mechanism to draw
global dependencies between input and output. The Transformer allows for
significantly more parallelization and can reach a new state of the art in
translation quality after being trained for as little as twelve hours on eight
P100 GPUs.

\section{Background}

The goal of reducing sequential computation also forms the foundation of the
Extended Neural GPU \cite{extendedneuralgpu}, ByteNet
\cite{kalchbrenner2016neural} and ConvS2S \cite{gehring2017convolutional}, all
of which use convolutional neural networks as basic building block, computing
hidden representations in parallel for all input and output positions. In
these models, the number of operations required to relate signals from two
arbitrary input or output positions grows in the distance between positions,
linearly for ConvS2S and logarithmically for ByteNet. This makes it more
difficult to learn dependencies between distant positions
\cite{hochreiter2001gradient}. In the Transformer this is reduced to a
constant number of operations, albeit at the cost of reduced effective
resolution due to averaging attention-weighted positions, an effect we
counteract with Multi-Head Attention as described in Section~3.2.

Self-attention, sometimes called intra-attention is an attention mechanism
relating different positions of a single sequence in order to compute a
representation of the sequence. Self-attention has been used successfully in a
variety of tasks including reading comprehension, abstractive summarization,
textual entailment and learning task-independent sentence representations
\cite{cheng2016long,decomposableattentionmodel,paulus2017deep,lin2017structured}.

End-to-end memory networks are based on a recurrent attention mechanism
instead of sequence-aligned recurrence and have been shown to perform well on
simple-language question answering and language modeling tasks
\cite{sukhbaatar2015}.

To the best of our knowledge, however, the Transformer is the first transduction
model relying entirely on self-attention to compute representations of its input
and output without using sequence-aligned RNNs or convolution. In the following
sections, we will describe the Transformer, motivate self-attention and discuss
its advantages over models such as
\cite{sutskever2014sequence,bahdanau2014neural} and
\cite{recurrentleastsquares,lecun2015deep}.

\section{Model Architecture}

Most competitive neural sequence transduction models have an encoder-decoder
structure \cite{cho2014learning,bahdanau2014neural,sutskever2014sequence}. Here,
the encoder maps an input sequence of symbol representations
$(x_1, \ldots, x_n)$ to a sequence of continuous representations
$\mathbf{z} = (z_1, \ldots, z_n)$. Given $\mathbf{z}$, the decoder then
generates an output sequence $(y_1, \ldots, y_m)$ of symbols one element at a
time. At each step the model is auto-regressive \cite{graves2013generating},
consuming the previously generated symbols as additional input when generating
the next.

The Transformer follows this overall architecture using stacked self-attention
and point-wise, fully connected layers for both the encoder and decoder, shown
in the left and right halves of Figure~1, respectively.

\subsection{Encoder and Decoder Stacks}

\textbf{Encoder:} The encoder is composed of a stack of $N = 6$ identical
layers. Each layer has two sub-layers. The first is a multi-head self-attention
mechanism, and the second is a simple, position-wise fully connected
feed-forward network. We employ a residual connection
\cite{he2016deep} around each of the two sub-layers, followed by layer
normalization \cite{ba2016layer}. That is, the output of each sub-layer is
$\mathrm{LayerNorm}(x + \mathrm{Sublayer}(x))$, where $\mathrm{Sublayer}(x)$
is the function implemented by the sub-layer itself. To facilitate these
residual connections, all sub-layers in the model, as well as the embedding
layers, produce outputs of dimension $d_{\text{model}} = 512$.

\textbf{Decoder:} The decoder is also composed of a stack of $N = 6$ identical
layers. In addition to the two sub-layers in each encoder layer, the decoder
inserts a third sub-layer, which performs multi-head attention over the output
of the encoder stack. Similar to the encoder, we employ residual connections
around each of the sub-layers, followed by layer normalization. We also modify
the self-attention sub-layer in the decoder stack to prevent positions from
attending to subsequent positions. This masking, combined with fact that the
output embeddings are offset by one position, ensures that the predictions for
position $i$ can depend only on the known outputs at positions less than $i$.

\subsection{Attention}

An attention function can be described as mapping a query and a set of
key-value pairs to an output, where the query, keys, values, and output are
all vectors. The output is computed as a weighted sum of the values, where the
weight assigned to each value is computed by a compatibility function of the
query with the corresponding key.

\subsubsection{Scaled Dot-Product Attention}

We call our particular attention ``Scaled Dot-Product Attention'' (Figure~2).
The input consists of queries and keys of dimension $d_k$, and values of
dimension $d_v$. We compute the dot products of the query with all keys, divide
each by $\sqrt{d_k}$, and apply a softmax function to obtain the weights on
the values.

In practice, we compute the attention function on a set of queries
simultaneously, packed together into a matrix $Q$. The keys and values are also
packed together into matrices $K$ and $V$. We compute the matrix of outputs as:
\[
  \mathrm{Attention}(Q, K, V) = \mathrm{softmax}\!\left(\frac{QK^{\top}}{\sqrt{d_k}}\right)V
\]

The two most commonly used attention functions are additive attention
\cite{bahdanau2014neural}, and dot-product (multiplicative) attention.
Dot-product attention is identical to our algorithm, except for the scaling
factor of $\frac{1}{\sqrt{d_k}}$. Additive attention computes the compatibility
function using a feed-forward network with a single hidden layer. While the two
are similar in theoretical complexity, dot-product attention is much faster and
more space-efficient in practice, since it can be implemented using highly
optimized matrix multiplication code.

While for small values of $d_k$ the two mechanisms perform similarly, additive
attention outperforms dot product attention without scaling for larger values of
$d_k$ \cite{DBLP:journals/corr/BritzGLL17}. We suspect that for large values
of $d_k$, the dot products grow large in magnitude, pushing the softmax function
into regions where it has extremely small gradients. To counteract this effect,
we scale the dot products by $\frac{1}{\sqrt{d_k}}$.

\subsubsection{Multi-Head Attention}

Instead of performing a single attention function with $d_{\text{model}}$-dimensional
keys, values and queries, we found it beneficial to linearly project the queries,
keys and values $h$ times with different, learned linear projections to $d_k$,
$d_k$ and $d_v$ dimensions, respectively. On each of these projected versions of
queries, keys and values we then perform the attention function in parallel,
yielding $d_v$-dimensional output values. These are concatenated and once again
projected, resulting in the final values, as depicted in Figure~2.

Multi-head attention allows the model to jointly attend to information from
different representation subspaces at different positions. With a single
attention head, averaging inhibits this.
\[
  \mathrm{MultiHead}(Q, K, V) = \mathrm{Concat}(\mathrm{head}_1, \ldots, \mathrm{head}_h)\,W^O
\]
where $\mathrm{head}_i = \mathrm{Attention}(QW_i^Q,\, KW_i^K,\, VW_i^V)$

Where the projections are parameter matrices
$W_i^Q \in \mathbb{R}^{d_{\text{model}} \times d_k}$,
$W_i^K \in \mathbb{R}^{d_{\text{model}} \times d_k}$,
$W_i^V \in \mathbb{R}^{d_{\text{model}} \times d_v}$ and
$W^O \in \mathbb{R}^{hd_v \times d_{\text{model}}}$.

In this work we employ $h = 8$ parallel attention layers, or heads. For each
of these we use $d_k = d_v = d_{\text{model}}/h = 64$. Due to the reduced
dimension of each head, the total computational cost is similar to that of
single-head attention with full dimensionality.

\subsubsection{Applications of Attention in our Model}

The Transformer uses multi-head attention in three different ways:

\begin{itemize}
  \item In ``encoder-decoder attention'' layers, the queries come from the
    previous decoder layer, and the memory keys and values come from the output
    of the encoder. This allows every position in the decoder to attend over all
    positions in the input sequence. This mimics the typical encoder-decoder
    attention mechanisms in sequence-to-sequence models such as
    \cite{wu2016google,bahdanau2014neural,decomposableattentionmodel}.
  \item The encoder contains self-attention layers. In a self-attention layer
    all of the keys, values and queries come from the same place, in this case,
    the output of the previous layer in the encoder. Each position in the
    encoder can attend to all positions in the previous layer of the encoder.
  \item Similarly, self-attention layers in the decoder allow each position in
    the decoder to attend to all positions in the decoder up to and including
    that position. We need to prevent leftward information flow in the decoder
    to preserve the auto-regressive property. We implement this inside of
    scaled dot-product attention by masking out (setting to $-\infty$) all
    values in the input to the softmax which correspond to illegal connections.
\end{itemize}

\subsection{Position-wise Feed-Forward Networks}

In addition to attention sub-layers, each of the layers in our encoder and
decoder contains a fully connected feed-forward network, which is applied to
each position separately and identically. This consists of two linear
transformations with a ReLU activation in between.
\[
  \mathrm{FFN}(x) = \max(0,\, xW_1 + b_1)\,W_2 + b_2
\]
While the linear transformations are the same across different positions, they
use different parameters from layer to layer. Another way of describing this is
as two convolutions with kernel size 1. The dimensionality of input and output
is $d_{\text{model}} = 512$, and the inner-layer has dimensionality
$d_{ff} = 2048$.

\subsection{Embeddings and Softmax}

Similarly to other sequence transduction models, we use learned embeddings to
convert the input tokens and output tokens to vectors of dimension
$d_{\text{model}}$. We also use the usual learned linear transformation and
softmax function to convert the decoder output to predicted next-token
probabilities. In our model, we share the same weight matrix between the two
embedding layers and the pre-softmax linear transformation, similar to
\cite{press2016using}. In the embedding layers, we multiply those weights by
$\sqrt{d_{\text{model}}}$.

\subsection{Positional Encoding}

Since our model contains no recurrence and no convolution, in order for the
model to make use of the order of the sequence, we must inject some information
about the relative or absolute position of the tokens in the sequence. To this
end, we add ``positional encodings'' to the input embeddings at the bottoms of
the encoder and decoder stacks. The positional encodings have the same dimension
$d_{\text{model}}$ as the embeddings, so that the two can be summed. There are
many choices of positional encodings, learned and fixed \cite{gehring2017convolutional}.

In this work, we use sine and cosine functions of different frequencies:
\[
  PE_{(pos,\,2i)}   = \sin\!\left(\frac{pos}{10000^{2i/d_{\text{model}}}}\right)
  \qquad
  PE_{(pos,\,2i+1)} = \cos\!\left(\frac{pos}{10000^{2i/d_{\text{model}}}}\right)
\]
where $pos$ is the position and $i$ is the dimension. That is, each dimension
of the positional encoding corresponds to a sinusoid. The wavelengths form a
geometric progression from $2\pi$ to $10000 \cdot 2\pi$. We chose this function
because we hypothesized it would allow the model to easily learn to attend by
relative positions, since for any fixed offset $k$,
$PE_{pos+k}$ can be represented as a linear function of $PE_{pos}$.

We also experimented with using learned positional embeddings
\cite{gehring2017convolutional} instead, and found that the two versions
produced nearly identical results (see Table~3 row (E)). We chose the sinusoidal
version because it may allow the model to extrapolate to sequence lengths longer
than the ones encountered during training.

\section{Why Self-Attention}

In this section we compare various aspects of self-attention layers to the
recurrent and convolutional layers commonly used for mapping one variable-length
sequence of symbol representations $(x_1, \ldots, x_n)$ to another sequence of
equal length $(z_1, \ldots, z_n)$, with $x_i, z_i \in \mathbb{R}^d$, such as a
hidden layer in a typical sequence transduction encoder or decoder. Motivating
our use of self-attention we consider three desiderata.

One is the total computational complexity per layer. Another is the amount of
computation that can be parallelized, as measured by the minimum number of
sequential operations required.

The third is the path length between long-range dependencies in the network.
Learning long-range dependencies is a key challenge in many sequence
transduction tasks. One key factor affecting the ability to learn such
dependencies is the length of the paths forward and backward signals have to
traverse in the network. The shorter these paths between any combination of
positions in the input and output sequences, the easier it is to learn
long-range dependencies \cite{hochreiter2001gradient}. Hence we also compare
the maximum path length between any two input and output positions in networks
composed of the different layer types.

As noted in Table~1, a self-attention layer connects all positions with a
constant number of sequentially executed operations, whereas a recurrent layer
requires $O(n)$ sequential operations. In terms of computational complexity,
self-attention layers are faster than recurrent layers when the sequence length
$n$ is smaller than the representation dimensionality $d$, which is most often
the case with sentence representations used by state-of-the-art models in
machine translations, such as word-piece \cite{wu2016google} and byte-pair
\cite{sennrich2015neural} representations.

To improve computational performance for tasks involving very long sequences,
self-attention could be restricted to considering only a neighborhood of size
$r$ in the input sequence centered around the respective output position. This
would increase the maximum path length to $O(n/r)$. We plan to investigate this
approach further in future work.

A single convolutional layer with kernel width $k < n$ does not connect all
pairs of input and output positions. Doing so requires a stack of $O(n/k)$
convolutional layers in the case of contiguous kernels, or $O(\log_k(n))$ in the
case of dilated convolutions \cite{yu2015multi}, increasing the length of the
longest paths between any two positions in the network. Convolutional layers are
generally more expensive than recurrent layers, by a factor of $k$. Separable
convolutions \cite{chollet2016xception}, however, decrease the complexity
considerably, to $O(k \cdot n \cdot d + n \cdot d^2)$. Even with $k = n$,
however, the complexity of a separable convolution is equal to the combination
of a self-attention layer and a point-wise feed-forward layer, the approach we
take in our model.

As side benefit, self-attention could yield more interpretable models. We
inspect attention distributions from our models and present and discuss examples
in the appendix. Not only do individual attention heads clearly learn to perform
different tasks, many appear to exhibit behavior related to the syntactic and
semantic structure of the sentences.

\section{Training}

This section describes the training regime for our models.

\subsection{Training Data and Batching}

We trained on the standard WMT 2014 English-German dataset consisting of about
4.5 million sentence pairs. Sentences were encoded using byte-pair encoding
\cite{sennrich2015neural}, which has a shared source-target vocabulary of about
37000 tokens. For English-French, we used the significantly larger WMT 2014
English-French dataset consisting of 36M sentences and split tokens into a
32000 word-piece vocabulary \cite{wu2016google}. Sentence pairs were batched
together by approximate sequence length. Each training batch contained a set of
sentence pairs containing approximately 25000 source tokens and 25000 target
tokens.

\subsection{Hardware and Schedule}

We trained our models on one machine with 8 NVIDIA P100 GPUs. For our base
models using the hyperparameters described throughout the paper, each training
step took about 0.4 seconds. We trained the base models for a total of 100,000
steps or 12 hours. For our big models, step time was 1.0 seconds. The big
models were trained for 300,000 steps (3.5 days).

\subsection{Optimizer}

We used the Adam optimizer \cite{kingma2014adam} with $\beta_1 = 0.9$,
$\beta_2 = 0.98$ and $\epsilon = 10^{-9}$. We varied the learning rate over the
course of training, according to the formula:
\[
  lrate = d_{\text{model}}^{-0.5} \cdot
  \min\!\left(step\_num^{-0.5},\;
  step\_num \cdot warmup\_steps^{-1.5}\right)
\]
This corresponds to increasing the learning rate linearly for the first
$warmup\_steps$ training steps, and decreasing it thereafter proportionally to
the inverse square root of the step number. We used $warmup\_steps = 4000$.

\subsection{Regularization}

We employ three types of regularization during training:

\textbf{Residual Dropout.} We apply dropout \cite{srivastava2014dropout} to the
output of each sub-layer, before it is added to the sub-layer input and
normalized. In addition, we apply dropout to the sums of the embeddings and the
positional encodings in both the encoder and decoder stacks. For the base model,
we use a rate of $P_{drop} = 0.1$.

\textbf{Label Smoothing.} During training, we employed label smoothing of value
$\epsilon_{ls} = 0.1$ \cite{szegedy2015rethinking}. This hurts perplexity, as
the model learns to be more unsure, but improves accuracy and BLEU score.

\section{Results}

\subsection{Machine Translation}

On the WMT 2014 English-to-German translation task, the big transformer model
(Transformer (big) in Table~2) outperforms the best previously reported models
including ensembles by more than 2.0 BLEU, establishing a new
state-of-the-art BLEU score of 28.4. The configuration of this model is listed
in the bottom line of Table~3. Training took 3.5 days on 8 P100 GPUs. Even our
base model surpasses all previously published models and ensembles, at a
fraction of the training cost of any of the competitive models.

On the WMT 2014 English-to-French translation task, our big model achieves a
BLEU score of 41.0, outperforming all of the previously published single models,
at less than $1/4$ the training cost of the previous state-of-the-art model.
The Transformer (big) model trained for English-to-French used dropout rate
$P_{drop} = 0.1$, instead of 0.3.

For the base models, we used a single model obtained by averaging the last 5
checkpoints, which were written at 10-minute intervals. For the big models, we
averaged the last 20 checkpoints. We used beam search with a beam size of 4 and
length penalty $\alpha = 0.6$ \cite{wu2016google}. These hyperparameters were
chosen after experimentation on the development set. We set the maximum output
length during inference to input length $+ 50$, but terminate early when
possible \cite{wu2016google}.

Table~2 summarizes our results and compares our translation quality and training
costs to other model architectures from the literature. We estimate the number
of floating point operations used to train a model by multiplying the training
time, the number of GPUs used, and an estimate of the sustained single-precision
floating-point capacity of each GPU.

\subsection{Model Variations}

To evaluate the importance of different components of the Transformer, we varied
our base model in different ways, measuring the change in performance on
English-to-German translation on the development set, newstest2013. We used beam
search as described in the previous section, but no checkpoint averaging. We
present these results in Table~3.

In Table~3 rows (A), we vary the number of attention heads and the attention key
and value dimensions, keeping the amount of computation constant, as described
in Section~3.2.2. While single-head attention is 0.9 BLEU worse than the best
setting, quality also drops off with too many heads.

In Table~3 rows (B), we observe that reducing the attention key size $d_k$ hurts
model quality. This suggests that determining compatibility is not easy and that
a more sophisticated compatibility function than dot product may be beneficial.
We further observe in rows (C) and (D) that, as expected, bigger models are
better, and dropout is very helpful in avoiding over-fitting. In row (E) we
replace our sinusoidal positional encoding with learned positional embeddings
\cite{gehring2017convolutional}, and observe nearly identical results to the
base model.

\subsection{English Constituency Parsing}

To evaluate if the Transformer can generalize to other tasks we performed
experiments on English constituency parsing. This task presents specific
challenges: the output is subject to strong structural constraints and is
significantly longer than the input. Furthermore, RNN sequence-to-sequence
models have not been able to attain state-of-the-art results in small-data
regimes \cite{vinyals2015grammar}.

We trained a 4-layer transformer with $d_{\text{model}} = 1024$ on the Wall
Street Journal (WSJ) portion of the Penn Treebank \cite{marcus1993building},
about 40K training sentences. We also trained it in a semi-supervised setting,
using the larger high-confidence and BerkleyParser corpora from with approximately
17M sentences \cite{vinyals2015grammar,zhu2013fast}. We used a vocabulary of
16K tokens for the WSJ only setting and a vocabulary of 32K tokens for the
semi-supervised setting.

We performed only a small number of experiments to select the dropout, both
attention and residual (section~5.4), learning rates and beam size on the
Section 22 development set, all other parameters remained unchanged from the
English-to-German base translation model. During inference, we increased the
maximum output length to input length $+ 300$. We used a beam size of 21 and
$\alpha = 0.3$ for both WSJ only and the semi-supervised setting.

Our results in Table~4 show that despite the lack of task-specific tuning our
model performs surprisingly well, yielding better results than all previously
reported models with the exception of the Recurrent Neural Network Grammar
\cite{dyer2016recurrent}.

In contrast to RNN sequence-to-sequence models \cite{vinyals2015grammar}, the
Transformer outperforms the BerkeleyParser \cite{petrov2006learning} even when
training only on the WSJ training set of 40K sentences.

\section{Conclusion}

In this work, we presented the Transformer, the first sequence transduction model
based entirely on attention, replacing the recurrent layers most commonly used in
encoder-decoder architectures with multi-headed self-attention.

For translation tasks, the Transformer can be trained significantly faster than
architectures based on recurrent or convolutional layers. We achieved a new state
of the art on both WMT 2014 English-to-German and WMT 2014 English-to-French
translation tasks. In the former task our best model outperforms even all
previously reported ensembles.

We are excited about the future of attention-based models and plan to apply them
to other tasks. We plan to extend the Transformer to problems involving input and
output modalities other than text and to investigate local, restricted attention
mechanisms to efficiently handle large inputs and outputs such as images, audio
and video. Making generation less sequential is another research goal of ours.

The code we used to train and evaluate our models is available at
\url{https://github.com/tensorflow/tensor2tensor}.

\chapter{《Deep Neural Networks for YouTube Recommendations》中文版}

\section{摘要}

YouTube 是现存规模最大、最复杂的工业级推荐系统之一。本文从高层视角描述该系统，重点介绍深度学习所带来的显著性能提升。全文按照经典的两阶段信息检索范式组织：首先详细介绍一个深度候选生成模型，然后描述一个独立的深度排序模型。我们还将分享在设计、迭代和维护这一对用户产生巨大影响的大规模推荐系统过程中所积累的实用经验与洞见。

\textbf{关键词：}推荐系统；深度学习；可扩展性

\section{引言}

YouTube 是全球最大的视频创作、分享与发现平台。YouTube 推荐系统负责帮助超过十亿用户从持续增长的视频库中发现个性化内容。本文重点关注深度学习近期对 YouTube 视频推荐系统所产生的巨大影响。图 1 展示了 YouTube 移动应用首页上的推荐界面。

YouTube 视频推荐面临以下三大主要挑战：

\begin{itemize}
  \item \textbf{规模（Scale）：}许多在小规模问题上表现良好的推荐算法在我们的规模下无法正常工作。处理 YouTube 庞大的用户群体和内容库，需要高度专业化的分布式学习算法和高效的服务系统。
  \item \textbf{时效性（Freshness）：}YouTube 的内容库极为动态，每秒都有大量视频上传。推荐系统需要足够灵敏，既能为新上传的内容建模，又能捕捉用户的最新行为。在探索/利用（exploration/exploitation）视角下，新内容与成熟视频之间的平衡是一个核心问题。
  \item \textbf{噪声（Noise）：}YouTube 用户的历史行为由于稀疏性及各种不可观测的外部因素，天然难以预测。我们几乎无法获得用户满意度的真实标签，只能对含噪的隐式反馈信号建模。此外，内容相关的元数据结构混乱，缺乏明确的本体定义。算法需要对训练数据的这些特性保持鲁棒性。
\end{itemize}

与谷歌旗下其他产品领域协同，YouTube 已完成向深度学习的根本性范式转变，将其作为几乎所有学习问题的通用解决方案。我们的系统构建在 Google Brain 之上，该平台近期以 TensorFlow 的形式开源。TensorFlow 提供了一个灵活的框架，可在大规模分布式训练环境下实验各种深度神经网络架构。我们的模型学习约十亿个参数，并在数千亿个样本上进行训练。

与大量矩阵分解方法的研究相比，将深度神经网络应用于推荐系统的工作相对较少。神经网络已被用于新闻推荐、引文推荐和评分预测；协同过滤被构建为深度神经网络，自编码器亦被用于此目的。Elkahky 等人利用深度学习进行跨域用户建模；在基于内容的场景中，Burges 等人将深度神经网络用于音乐推荐。

本文的组织结构如下：第 2 节简要介绍系统概述；第 3 节更详细地描述候选生成模型，包括其训练方式以及如何用于服务推荐，实验结果将展示该模型如何从更深的隐藏层和额外的异质信号中受益；第 4 节详细介绍排序模型，包括如何改进经典逻辑回归以训练预测期望观看时长（而非点击概率）的模型，实验结果将表明隐藏层深度在此同样有益；最后，第 5 节呈现我们的结论与经验教训。

\section{系统概述}

我们推荐系统的整体架构如图 2 所示。该系统由两个神经网络组成：一个用于\textit{候选生成（candidate generation）}，另一个用于\textit{排序（ranking）}。

候选生成网络以用户在 YouTube 上的活动历史作为输入，从庞大的视频库中检索出一小部分（数百个）可能与用户相关的视频。这些候选视频具有较高的精确率，能够总体上满足用户的相关性需求。候选生成网络仅通过协同过滤提供粗粒度的个性化。用户之间的相似性通过粗粒度特征来表达，例如观看视频的 ID、搜索词令牌以及人口统计信息。

在列表中呈现少量"最佳"推荐，需要对候选视频的相对重要性进行细粒度的区分，以实现高召回率。排序网络通过使用描述视频与用户的丰富特征集，依据目标函数为每个视频分配一个分数来完成这一任务。得分最高的视频按分数排序后呈现给用户。

两阶段推荐方法使我们能够从数百万规模的庞大视频库中生成推荐，同时确保最终呈现在用户设备上的少量视频是个性化的、能够吸引用户的。此外，这一设计还支持融合来自其他来源的候选视频，例如早期工作中描述的那些来源。

在开发过程中，我们大量使用离线指标（精确率、召回率、排序损失等）来指导系统的迭代改进。然而，对于算法或模型有效性的最终判断，我们依赖于在线 A/B 实验。在线实验可以衡量点击率、观看时长以及其他用户参与度指标的细微变化。这一点至关重要，因为在线 A/B 实验结果并不总是与离线实验结果相关。

\section{候选生成}

在候选生成阶段，庞大的 YouTube 视频库被精简为数百个可能与用户相关的视频。本文所描述的推荐系统的前身是一种在排序损失下训练的矩阵分解方法。我们神经网络模型的早期迭代版本通过仅嵌入用户历史观看记录的浅层网络来模拟这种分解行为。从这个角度来看，我们的方法可以视为分解技术的非线性推广。

\subsection{将推荐建模为分类}

我们将推荐问题建模为极端多分类问题，预测任务转化为：给定用户 $U$ 和上下文 $C$，在包含数百万视频 $i$（类别）的视频库 $V$ 中，精确分类用户在时刻 $t$ 将观看的特定视频 $w_t$，
\[
  P(w_t = i \mid U, C) = \frac{e^{v_i u}}{\sum_{j \in V} e^{v_j u}}
\]
其中 $u \in \mathbb{R}^N$ 表示用户与上下文对的高维"嵌入（embedding）"，$v_j \in \mathbb{R}^N$ 表示各候选视频的嵌入。在此设定下，嵌入即是将稀疏实体（单个视频、用户等）映射到 $\mathbb{R}^N$ 中稠密向量的映射。深度神经网络的任务是以用户的历史记录和上下文为输入，学习用户嵌入 $u$，使其能够通过 softmax 分类器对视频进行有效区分。

尽管 YouTube 上存在显式反馈机制（点赞/点踩、产品内调查等），我们仍选择使用观看行为这一隐式反馈来训练模型，以用户看完视频作为正样本。这一选择基于以下考量：可用的隐式用户历史数据在量级上远超显式反馈，使我们能够为长尾内容生成推荐，而在这些领域显式反馈极为稀疏。

\textit{高效的极端多分类（Efficient Extreme Multiclass）}

为了在数百万类别下高效训练此类模型，我们采用从背景分布中采样负类（"候选采样"）的技术，并通过重要性加权对采样偏差进行校正。对于每个样本，交叉熵损失在真实标签和采样的负类上同时被最小化。实践中采样数千个负样本，相比传统 softmax 实现了超过 100 倍的加速。层次化 softmax 是另一种常见替代方案，但我们未能达到可比的精度。在层次化 softmax 中，遍历树中每个节点需要在通常不相关的类别集合之间进行判别，这使分类问题更加困难，导致性能下降。

在服务阶段，需要计算最可能的 $N$ 个类别（视频）以选出排名前 $N$ 的视频呈现给用户。在数十毫秒的严格服务延迟下对数百万候选项打分，需要一种次线性复杂度的近似评分方案。YouTube 此前的系统依赖哈希方法，本文描述的分类器采用了类似的方案。由于服务时不需要 softmax 输出层的校准似然值，评分问题退化为点积空间中的最近邻搜索，可使用通用库实现。我们发现 A/B 实验结果对最近邻搜索算法的选择并不十分敏感。

\subsection{模型架构}

受连续词袋语言模型的启发，我们在固定词表中为每个视频学习高维嵌入，并将这些嵌入输入前馈神经网络。用户的观看历史由一个变长的稀疏视频 ID 序列表示，通过嵌入映射为稠密向量表示。网络需要固定大小的稠密输入，实验表明对嵌入取平均在多种策略（求和、逐元素最大值等）中表现最佳。重要的是，嵌入与其他所有模型参数通过正常的梯度下降反向传播联合学习。特征被拼接成一个宽的第一层，随后是若干全连接整流线性单元（ReLU）层。图 3 展示了加入以下所述非视频观看特征后的通用网络架构。

\subsection{异质信号}

将深度神经网络作为矩阵分解推广的一个关键优势在于，任意连续型和类别型特征都可以方便地加入模型。搜索历史的处理方式与观看历史类似——每个查询被切分为 unigram 和 bigram，每个 token 被嵌入。取平均后，用户经过分词和嵌入处理的查询序列代表了一份压缩的稠密搜索历史。人口统计特征对于为新用户提供先验信息、使推荐行为合理化至关重要。用户的地理区域和设备被嵌入后拼接。性别、登录状态和年龄等简单二值型和连续型特征以归一化到 $[0, 1]$ 的实数值直接输入网络。

\textit{"样本年龄"特征（"Example Age" Feature）}

每秒都有大量视频上传至 YouTube。推荐新上传的（"新鲜"）内容对 YouTube 作为产品而言极为重要。我们持续观察到用户偏好新鲜内容，但不以牺牲相关性为代价。除了简单推荐用户想看的新视频这一一阶效应外，还存在引导和传播病毒式内容这一关键的二阶现象。

机器学习系统通常对历史存在隐式偏差，因为它们是从历史样本中预测未来行为训练而来的。视频流行度的分布高度非平稳，但我们推荐系统所产生的语料库上的多项分布将反映数周训练窗口内的平均观看可能性。为了纠正这一偏差，我们在训练时将训练样本的年龄（age）作为特征输入。在服务时，该特征被设为零（或略为负值），以体现模型是在训练窗口末端做出预测。

图 4 展示了该方法在一个任意选取的视频上的有效性。

\subsection{标签与上下文选择}

需要特别强调的是，推荐往往涉及求解一个\textit{代理问题（surrogate problem）}并将结果迁移到特定场景。一个经典例子是"准确预测评分可以带来有效电影推荐"这一假设。我们发现，代理学习问题的选择对 A/B 测试中的性能影响极大，但在离线实验中却很难衡量。

训练样本来自所有 YouTube 观看记录（包括嵌入在其他网站上的观看），而非仅限于我们推荐产生的观看。否则，新内容将极难浮现，推荐系统也会对利用（exploitation）产生过度偏倚。若用户通过推荐以外的途径发现视频，我们希望能够通过协同过滤迅速将这一发现传播给其他用户。另一个改善在线指标的关键洞见是：为每位用户生成固定数量的训练样本，从而在损失函数中对用户进行等权重处理，防止少数高活跃用户主导损失。

有些违反直觉的是，为防止模型利用网站结构并对代理问题过拟合，必须谨慎地\textit{对分类器隐藏某些信息}。举例来说，若用户刚刚搜索了"taylor swift"，由于问题被建模为预测下一个观看视频，给定该信息的分类器会预测最可能被观看的视频正是对应搜索结果页上出现的视频。不出所料，将用户最近的搜索页面作为首页推荐的表现极差。通过丢弃序列信息、以无序词袋表示搜索查询，分类器不再能直接感知标签的来源。

视频的自然消费模式通常导致非常不对称的共同观看概率。系列剧通常按顺序观看，用户往往从某类型中最广为人知的内容开始，逐渐深入到更细分的领域。因此，我们发现预测用户的\textit{下一次}观看，比预测随机留出的观看效果更好（图 5）。许多协同过滤系统通过留出随机一项并从用户历史中的其他项进行预测来隐式选择标签和上下文（图 5a），这会泄露未来信息并忽略非对称的消费模式。与之相反，我们通过选择一次随机观看来"回滚"用户历史，仅将用户在留出标签观看之前的行为作为输入（图 5b）。

\subsection{特征与深度实验}

如图 6 所示，增加特征和深度能显著提升留出数据上的精确率。在这些实验中，100 万个视频词表和 100 万个搜索词令牌词表各以 256 个浮点数嵌入，最大词袋大小为 50 条最近观看和 50 条最近搜索。softmax 层在相同的 100 万视频类别上输出维度为 256 的多项分布（可视为一个独立的输出视频嵌入）。这些模型在所有 YouTube 用户上训练至收敛，对应在数据上进行若干个 epoch。网络结构遵循常见的"塔形"模式——网络底层最宽，每层隐藏层的单元数减半（类似图 3）。深度为零的网络实际上是一个线性分解方案，与前代系统的表现非常相似。不断增加宽度和深度，直到增量收益减弱、收敛变得困难为止：

\begin{itemize}
  \item 深度 0：一个线性层，将拼接层变换至 softmax 维度 256
  \item 深度 1：256 ReLU
  \item 深度 2：512 ReLU $\rightarrow$ 256 ReLU
  \item 深度 3：1024 ReLU $\rightarrow$ 512 ReLU $\rightarrow$ 256 ReLU
  \item 深度 4：2048 ReLU $\rightarrow$ 1024 ReLU $\rightarrow$ 512 ReLU $\rightarrow$ 256 ReLU
\end{itemize}

\section{排序}

排序的核心职责是利用曝光数据，针对特定用户界面对候选预测结果进行专项化校准。例如，某用户总体上以较高概率观看某视频，但由于缩略图的选择，未必会点击该视频在首页上的特定曝光位。在排序阶段，由于只需对数百个视频打分（而非候选生成阶段的数百万个），我们能够获取大量描述视频与用户关系的特征。排序对于融合来自不同候选来源、分数不可直接比较的结果同样至关重要。

我们使用与候选生成架构类似的深度神经网络，通过逻辑回归（图 7）为每个视频曝光独立打分，然后按分数排序返回给用户。最终排序目标依据在线 A/B 测试结果持续调整，但通常是每次曝光期望观看时长的简单函数。按点击率排序往往会推广用户不会看完的欺骗性视频（"标题党"），而观看时长能更好地衡量用户参与度。

\subsection{特征表示}

我们的特征遵循传统的类别型与连续型/有序型特征分类体系。所使用的类别型特征在基数上差异悬殊——有些是二值型（如用户是否登录），另一些则有数百万个可能取值（如用户的最近一次搜索词）。特征还按照是贡献单一值（"单值型"）还是一组值（"多值型"）进一步区分。单值类别型特征的典型例子是当前打分的视频 ID，对应的多值特征则可以是用户观看的最近 $N$ 个视频 ID 的词袋。此外，特征还按是否描述候选项属性（"曝光"特征）或用户/上下文属性（"查询"特征）进行分类。查询特征每次请求只计算一次，而曝光特征则需对每个打分候选项分别计算。

\textit{特征工程（Feature Engineering）}

我们的排序模型通常使用数百个特征，类别型与连续型特征大致各占一半。尽管深度学习有望减轻手工特征工程的负担，但我们的原始数据性质使其无法直接输入前馈神经网络。我们仍然投入大量工程资源将用户和视频数据转化为有用的特征，核心挑战在于如何表示用户行为的时序序列，以及这些行为与当前打分视频曝光之间的关系。

我们观察到最重要的信号，是描述用户与候选项本身及相似候选项的历史交互情况的特征，与其他排序广告经验相吻合。例如，考虑用户与当前打分视频所属频道的历史关系——用户观看了该频道多少视频？用户最近一次观看该主题内容是什么时候？这些描述用户在相关内容上历史行为的连续型特征尤为强大，因为它们能在不同内容之间良好泛化。我们还发现，将候选生成阶段的信息以特征形式传递给排序阶段至关重要，例如：哪些来源提名了这个候选视频？它们分别给出了怎样的分数？

描述历史视频曝光频率的特征对于在推荐中引入"扰动（churn）"同样关键——确保连续请求不会返回完全相同的列表。若某视频最近被推荐给用户但未被观看，模型将在下次页面加载时自然地降低该曝光的排名。实时提供最新的曝光和观看历史本身就是一项重大的工程挑战，超出了本文讨论范围，但对于提供响应灵敏的推荐至关重要。

\textit{嵌入类别型特征（Embedding Categorical Features）}

与候选生成类似，我们使用嵌入将稀疏类别型特征映射为适合神经网络的稠密表示。每个唯一 ID 空间（"词表"）有一个独立的可学习嵌入，其维度大致与唯一取值数量的对数成正比。这些词表是在训练开始前对数据进行一次遍历后构建的简单查找表。极高基数的 ID 空间（如视频 ID 或搜索词）通过仅保留按点击曝光频率排序的前 $N$ 个进行截断。词表外的取值简单地映射到零嵌入。与候选生成相同，多值类别型特征的嵌入在输入网络前取平均。

重要的是，同一 ID 空间内的类别型特征共享底层嵌入。例如，存在一个单一的全局视频 ID 嵌入，被许多不同特征使用（曝光的视频 ID、用户最近观看的视频 ID、"种子"推荐的视频 ID 等）。尽管共享嵌入，每个特征仍单独输入网络，使上层网络能够为每个特征学习专门的表示。共享嵌入有助于改善泛化、加速训练并降低内存需求。绝大多数模型参数都集中在这些高基数嵌入空间中——例如，100 万个 ID 在 32 维空间中嵌入的参数量是 2048 单元全连接层的 7 倍。

\textit{连续型特征归一化（Normalizing Continuous Features）}

神经网络对输入的尺度和分布极为敏感，而决策树集成等方法则对单个特征的缩放不变。我们发现，对连续型特征进行合理归一化对收敛至关重要。对于分布为 $f$ 的连续型特征 $x$，通过累积分布函数 $\tilde{x} = \int_{-\infty}^{x} df$ 将其缩放至在 $[0, 1)$ 上均匀分布，得到归一化特征 $\tilde{x}$。该积分在训练开始前对数据进行一次遍历，通过特征值分位数上的线性插值近似。

除原始归一化特征 $\tilde{x}$ 外，我们还将其幂次 $\tilde{x}^2$ 和 $\sqrt{\tilde{x}}$ 一并输入，赋予网络更强的表达能力，使其能够轻松构建特征的超线性和次线性函数。实验表明，输入连续型特征的幂次能够改善离线准确率。

\subsection{期望观看时长建模}

我们的目标是根据正样本（视频曝光被点击）和负样本（曝光未被点击）的训练样本，预测期望观看时长。正样本被标注了用户观看该视频所花费的时间。为预测期望观看时长，我们采用加权逻辑回归技术，该技术正是为此目的而开发的。

模型在交叉熵损失下以逻辑回归训练（图 7）。然而，正样本（被点击的曝光）以该视频的实际观看时长加权，负样本（未被点击的曝光）则统一赋予单位权重。由此，逻辑回归所学习的赔率为 $\frac{\sum T_i}{N - k}$，其中 $N$ 为训练样本总数，$k$ 为正样本数，$T_i$ 为第 $i$ 个样本的观看时长。假设正样本比例较小（在我们的场景下确实如此），学习到的赔率近似为 $E[T](1 + P)$，其中 $P$ 为点击概率，$E[T]$ 为该曝光的期望观看时长。由于 $P$ 较小，该乘积接近于 $E[T]$。在推断时，我们使用指数函数 $e^x$ 作为最终激活函数，以产生对期望观看时长的近似估计。

\subsection{隐藏层实验}

表 1 展示了我们在不同隐藏层配置下的次日留出数据结果。每种配置的指标值（"加权逐用户损失"）通过考察单次页面加载中呈现给用户的正样本（被点击）和负样本（未被点击）曝光对来计算。我们首先用模型对两个曝光打分：若负样本得分高于正样本，则认为正样本的观看时长被错误预测。加权逐用户损失即为留出曝光对中，被错误预测的观看时长占总观看时长的比例。

结果表明，增加隐藏层宽度和深度均能改善性能。然而，代价是推断所需的服务端 CPU 时间。$1024 \rightarrow 512 \rightarrow 256$ 的 ReLU 配置在满足服务 CPU 预算的前提下取得了最佳结果。

对于 $1024 \rightarrow 512 \rightarrow 256$ 模型，我们尝试仅输入归一化连续型特征而不加入其幂次，损失增加了 0.2\%。在相同隐藏层配置下，我们还训练了一个对正负样本等权重处理的模型，正如预期，观看时长加权损失显著上升了 4.1\%。

\section{结论}

我们描述了用于 YouTube 视频推荐的深度神经网络架构，将其拆分为两个独立问题：候选生成与排序。

我们的深度协同过滤模型能够有效融合多种信号，并通过多层深度对信号间的交互建模，优于 YouTube 此前使用的矩阵分解方法。在为推荐选择代理问题时，艺术性多于科学性——我们发现，预测用户的未来观看在在线指标上表现优异，这得益于它能捕捉非对称的共同观看行为，并防止未来信息的泄露。对分类器隐藏判别性信号同样是取得良好结果的关键——否则模型会对代理问题过拟合，无法良好迁移到首页场景。

我们证明，将训练样本的年龄作为输入特征，能够消除模型对历史的隐式偏倚，使模型能够表示视频流行度的时间依赖行为。这改善了离线留出精确率，并在 A/B 测试中显著提升了近期上传视频的观看时长。

排序是一个更为经典的机器学习问题，但我们的深度学习方法在观看时长预测上优于此前的线性和基于树的方法。推荐系统尤其受益于描述用户对候选项历史行为的专项特征。深度神经网络要求对类别型和连续型特征进行特殊表示，我们分别通过嵌入和分位数归一化进行处理。多层深度被证明能够有效建模数百个特征间的非线性交互。

逻辑回归通过对正样本以观看时长加权、对负样本赋予单位权重进行改进，使我们能够学习到对期望观看时长精确建模的赔率。该方法在以观看时长加权的排序评估指标上，远优于直接预测点击率的方法。

\chapter{Deep Neural Networks for YouTube Recommendations}

\section{Abstract}

YouTube represents one of the largest scale and most sophisticated industrial
recommendation systems in existence. In this paper, we describe the system at
a high level and focus on the dramatic performance improvements brought by
deep learning. The paper is split according to the classic two-stage
information retrieval dichotomy: first, we detail a deep candidate generation
model and then describe a separate deep ranking model. We also provide
practical lessons and insights derived from designing, iterating and
maintaining a massive recommendation system with enormous user-facing impact.

\textbf{Keywords:} recommender system; deep learning; scalability

\section{Introduction}

YouTube is the world's largest platform for creating, sharing and discovering
video content. YouTube recommendations are responsible for helping more than
a billion users discover personalized content from an ever-growing corpus of
videos. In this paper we will focus on the immense impact deep learning has
recently had on the YouTube video recommendations system. Figure~1
illustrates the recommendations on the YouTube mobile app home.

Recommending YouTube videos is extremely challenging from three major
perspectives:

\begin{itemize}
  \item \textbf{Scale:} Many existing recommendation algorithms proven to
    work well on small problems fail to operate on our scale. Highly
    specialized distributed learning algorithms and efficient serving systems
    are essential for handling YouTube's massive user base and corpus.
  \item \textbf{Freshness:} YouTube has a very dynamic corpus with many hours
    of video are uploaded per second. The recommendation system should be
    responsive enough to model newly uploaded content as well as the latest
    actions taken by the user. Balancing new content with well-established
    videos can be understood from an exploration/exploitation perspective.
  \item \textbf{Noise:} Historical user behavior on YouTube is inherently
    difficult to predict due to sparsity and a variety of unobservable
    external factors. We rarely obtain the ground truth of user satisfaction
    and instead model noisy implicit feedback signals. Furthermore, metadata
    associated with content is poorly structured without a well defined
    ontology. Our algorithms need to be robust to these particular
    characteristics of our training data.
\end{itemize}

In conjugation with other product areas across Google, YouTube has undergone
a fundamental paradigm shift towards using deep learning as a general-purpose
solution for nearly all learning problems. Our system is built on Google
Brain~\cite{dean2012large} which was recently open sourced as
TensorFlow~\cite{tensorflow2015}. TensorFlow provides a flexible framework
for experimenting with various deep neural network architectures using
large-scale distributed training. Our models learn approximately one billion
parameters and are trained on hundreds of billions of examples.

In contrast to vast amount of research in matrix factorization
methods~\cite{su2009survey}, there is relatively little work using deep
neural networks for recommendation systems. Neural networks are used for
recommending news in~\cite{oh2014personalized}, citations
in~\cite{huang2015neural} and review ratings
in~\cite{tang2015user}. Collaborative filtering is formulated as a deep
neural network in~\cite{wang2015collaborative} and autoencoders
in~\cite{sedhain2015autorec}. Elkahky et al.\ used deep learning for cross
domain user modeling~\cite{elkahky2015multi}. In a content-based setting,
Burges et al.\ used deep neural networks for music
recommendation~\cite{vandenoord2013deep}.

The paper is organized as follows: A brief system overview is presented in
Section~2. Section~3 describes the candidate generation model in more detail,
including how it is trained and used to serve recommendations. Experimental
results will show how the model benefits from deep layers of hidden units and
additional heterogeneous signals. Section~4 details the ranking model,
including how classic logistic regression is modified to train a model
predicting expected watch time (rather than click probability). Experimental
results will show that hidden layer depth is helpful as well in this
situation. Finally, Section~5 presents our conclusions and lessons learned.

\section{System Overview}

The overall structure of our recommendation system is illustrated in
Figure~2. The system is comprised of two neural networks: one for
\textit{candidate generation} and one for \textit{ranking}.

The candidate generation network takes events from the user's YouTube
activity history as input and retrieves a small subset (hundreds) of videos
from a large corpus. These candidates are intended to be generally relevant
to the user with high precision. The candidate generation network only
provides broad personalization via collaborative filtering. The similarity
between users is expressed in terms of coarse features such as IDs of video
watches, search query tokens and demographics.

Presenting a few ``best'' recommendations in a list requires a fine-level
representation to distinguish relative importance among candidates with high
recall. The ranking network accomplishes this task by assigning a score to
each video according to a desired objective function using a rich set of
features describing the video and user. The highest scoring videos are
presented to the user, ranked by their score.

The two-stage approach to recommendation allows us to make recommendations
from a very large corpus (millions) of videos while still being certain that
the small number of videos appearing on the device are personalized and
engaging for the user. Furthermore, this design enables blending candidates
generated by other sources, such as those described in an earlier
work~\cite{davidson2010youtube}.

During development, we make extensive use of offline metrics (precision,
recall, ranking loss, etc.) to guide iterative improvements to our system.
However for the final determination of the effectiveness of an algorithm or
model, we rely on A/B testing via live experiments. In a live experiment, we
can measure subtle changes in click-through rate, watch time, and many other
metrics that measure user engagement. This is important because live A/B
results are not always correlated with offline experiments.

\section{Candidate Generation}

During candidate generation, the enormous YouTube corpus is winnowed down to
hundreds of videos that may be relevant to the user. The predecessor to the
recommender described here was a matrix factorization approach trained under
rank loss~\cite{weston2011wsabie}. Early iterations of our neural network
model mimicked this factorization behavior with shallow networks that only
embedded the user's previous watches. From this perspective, our approach
can be viewed as a non-linear generalization of factorization techniques.

\subsection{Recommendation as Classification}

We pose recommendation as extreme multiclass classification where the
prediction problem becomes accurately classifying a specific video watch
$w_t$ at time $t$ among millions of videos $i$ (classes) from a corpus $V$
based on a user $U$ and context $C$,
\[
  P(w_t = i \mid U, C) = \frac{e^{v_i u}}{\sum_{j \in V} e^{v_j u}}
\]
where $u \in \mathbb{R}^N$ represents a high-dimensional ``embedding'' of the
user, context pair and the $v_j \in \mathbb{R}^N$ represent embeddings of
each candidate video. In this setting, an embedding is simply a mapping of
sparse entities (individual videos, users etc.) into a dense vector in
$\mathbb{R}^N$. The task of the deep neural network is to learn user
embeddings $u$ as a function of the user's history and context that are
useful for discriminating among videos with a softmax classifier.

Although explicit feedback mechanisms exist on YouTube (thumbs up/down,
in-product surveys, etc.) we use the implicit feedback~\cite{oard1998implicit}
of watches to train the model, where a user completing a video is a positive
example. This choice is based on the orders of magnitude more implicit user
history available, allowing us to produce recommendations deep in the tail
where explicit feedback is extremely sparse.

\textit{Efficient Extreme Multiclass}

To efficiently train such a model with millions of classes, we rely on a
technique to sample negative classes from the background distribution
(``candidate sampling'') and then correct for this sampling via importance
weighting~\cite{jean2014using}. For each example the cross-entropy loss is
minimized for the true label and the sampled negative classes. In practice
several thousand negatives are sampled, corresponding to more than 100 times
speedup over traditional softmax. A popular alternative approach is
hierarchical softmax~\cite{morin2005hierarchical}, but we weren't able to
achieve comparable accuracy. In hierarchical softmax, traversing each node
in the tree involves discriminating between sets of classes that are often
unrelated, making the classification problem much more difficult and
degrading performance.

At serving time we need to compute the most likely $N$ classes (videos) in
order to choose the top $N$ to present to the user. Scoring millions of items
under a strict serving latency of tens of milliseconds requires an
approximate scoring scheme sublinear in the number of classes. Previous
systems at YouTube relied on hashing~\cite{weston2013label} and the
classifier described here uses a similar approach. Since calibrated
likelihoods from the softmax output layer are not needed at serving time, the
scoring problem reduces to a nearest neighbor search in the dot product space
for which general purpose libraries can be used~\cite{liu2004investigation}.
We found that A/B results were not particularly sensitive to the choice of
nearest neighbor search algorithm.

\subsection{Model Architecture}

Inspired by continuous bag of words language models~\cite{mikolov2013distributed},
we learn high dimensional embeddings for each video in a fixed vocabulary and
feed these embeddings into a feedforward neural network. A user's watch
history is represented by a variable-length sequence of sparse video IDs
which is mapped to a dense vector representation via the embeddings. The
network requires fixed-sized dense inputs and simply averaging the embeddings
performed best among several strategies (sum, component-wise max, etc.).
Importantly, the embeddings are learned jointly with all other model
parameters through normal gradient descent backpropagation updates. Features
are concatenated into a wide first layer, followed by several layers of fully
connected Rectified Linear Units (ReLU)~\cite{glorot2011deep}. Figure~3 shows
the general network architecture with additional non-video watch features
described below.

\subsection{Heterogeneous Signals}

A key advantage of using deep neural networks as a generalization of matrix
factorization is that arbitrary continuous and categorical features can be
easily added to the model. Search history is treated similarly to watch
history --- each query is tokenized into unigrams and bigrams and each token
is embedded. Once averaged, the user's tokenized, embedded queries represent
a summarized dense search history. Demographic features are important for
providing priors so that the recommendations behave reasonably for new users.
The user's geographic region and device are embedded and concatenated. Simple
binary and continuous features such as the user's gender, logged-in state and
age are input directly into the network as real values normalized to
$[0, 1]$.

\textit{``Example Age'' Feature}

Many hours worth of videos are uploaded each second to YouTube. Recommending
this recently uploaded (``fresh'') content is extremely important for
YouTube as a product. We consistently observe that users prefer fresh
content, though not at the expense of relevance. In addition to the
first-order effect of simply recommending new videos that users want to
watch, there is a critical secondary phenomenon of bootstrapping and
propagating viral content~\cite{jiang2014viral}.

Machine learning systems often exhibit an implicit bias towards the past
because they are trained to predict future behavior from historical examples.
The distribution of video popularity is highly non-stationary but the
multinomial distribution over the corpus produced by our recommender will
reflect the average watch likelihood in the training window of several weeks.
To correct for this, we feed the age of the training example as a feature
during training. At serving time, this feature is set to zero (or slightly
negative) to reflect that the model is making predictions at the very end of
the training window.

Figure~4 demonstrates the efficacy of this approach on an arbitrarily chosen
video~\cite{zayn2016pillowtalk}.

\subsection{Label and Context Selection}

It is important to emphasize that recommendation often involves solving a
\textit{surrogate problem} and transferring the result to a particular
context. A classic example is the assumption that accurately predicting
ratings leads to effective movie recommendations~\cite{amatriain2012building}.
We have found that the choice of this surrogate learning problem has an
outsized importance on performance in A/B testing but is very difficult to
measure with offline experiments.

Training examples are generated from all YouTube watches (even those
embedded on other sites) rather than just watches on the recommendations we
produce. Otherwise, it would be very difficult for new content to surface and
the recommender would be overly biased towards exploitation. If users are
discovering videos through means other than our recommendations, we want to
be able to quickly propagate this discovery to others via collaborative
filtering. Another key insight that improved live metrics was to generate a
fixed number of training examples per user, effectively weighting our users
equally in the loss function. This prevented a small cohort of highly active
users from dominating the loss.

Somewhat counter-intuitively, great care must be taken to \textit{withhold
information from the classifier} in order to prevent the model from
exploiting the structure of the site and overfitting the surrogate problem.
Consider as an example a case in which the user has just issued a search
query for ``taylor swift''. Since our problem is posed as predicting the next
watched video, a classifier given this information will predict that the
most likely videos to be watched are those which appear on the corresponding
search results page for ``taylor swift''. Unsurpisingly, reproducing the
user's last search page as homepage recommendations performs very poorly. By
discarding sequence information and representing search queries with an
unordered bag of tokens, the classifier is no longer directly aware of the
origin of the label.

Natural consumption patterns of videos typically lead to very asymmetric
co-watch probabilities. Episodic series are usually watched sequentially and
users often discover artists in a genre beginning with the most broadly
popular before focusing on smaller niches. We therefore found much better
performance predicting the user's next watch, rather than predicting a
randomly held-out watch (Figure~5). Many collaborative filtering systems
implicitly choose the labels and context by holding out a random item and
predicting it from other items in the user's history (5a). This leaks future
information and ignores any asymmetric consumption patterns. In contrast, we
``rollback'' a user's history by choosing a random watch and only input
actions the user took before the held-out label watch (5b).

\subsection{Experiments with Features and Depth}

Adding features and depth significantly improves precision on holdout data
as shown in Figure~6. In these experiments, a vocabulary of 1M videos and 1M
search tokens were embedded with 256 floats each in a maximum bag size of 50
recent watches and 50 recent searches. The softmax layer outputs a
multinomial distribution over the same 1M video classes with a dimension of
256 (which can be thought of as a separate output video embedding). These
models were trained until convergence over all YouTube users, corresponding
to several epochs over the data. Network structure followed a common
``tower'' pattern in which the bottom of the network is widest and each
successive hidden layer halves the number of units (similar to Figure~3).
The depth zero network is effectively a linear factorization scheme which
performed very similarly to the predecessor system. Width and depth were
added until the incremental benefit diminished and convergence became
difficult:

\begin{itemize}
  \item Depth 0: A linear layer simply transforms the concatenation layer to
    match the softmax dimension of 256
  \item Depth 1: 256 ReLU
  \item Depth 2: 512 ReLU $\rightarrow$ 256 ReLU
  \item Depth 3: 1024 ReLU $\rightarrow$ 512 ReLU $\rightarrow$ 256 ReLU
  \item Depth 4: 2048 ReLU $\rightarrow$ 1024 ReLU $\rightarrow$ 512 ReLU
    $\rightarrow$ 256 ReLU
\end{itemize}

\section{Ranking}

The primary role of ranking is to use impression data to specialize and
calibrate candidate predictions for the particular user interface. For
example, a user may watch a given video with high probability generally but
is unlikely to click on the specific homepage impression due to the choice
of thumbnail image. During ranking, we have access to many more features
describing the video and the user's relationship to the video because only a
few hundred videos are being scored rather than the millions scored in
candidate generation. Ranking is also crucial for ensembling different
candidate sources whose scores are not directly comparable.

We use a deep neural network with similar architecture as candidate
generation to assign an independent score to each video impression using
logistic regression (Figure~7). The list of videos is then sorted by this
score and returned to the user. Our final ranking objective is constantly
being tuned based on live A/B testing results but is generally a simple
function of expected watch time per impression. Ranking by click-through
rate often promotes deceptive videos that the user does not complete
(``clickbait'') whereas watch time better captures
engagement~\cite{meyerson2012watchtime,yi2014beyond}.

\subsection{Feature Representation}

Our features are segregated with the traditional taxonomy of categorical and
continuous/ordinal features. The categorical features we use vary widely in
their cardinality --- some are binary (e.g.\ whether the user is logged-in)
while others have millions of possible values (e.g.\ the user's last search
query). Features are further split according to whether they contribute only
a single value (``univalent'') or a set of values (``multivalent''). An
example of a univalent categorical feature is the video ID of the impression
being scored, while a corresponding multivalent feature might be a bag of the
last $N$ video IDs the user has watched. We also classify features according
to whether they describe properties of the item (``impression'') or
properties of the user/context (``query''). Query features are computed once
per request while impression features are computed for each item scored.

\textit{Feature Engineering}

We typically use hundreds of features in our ranking models, roughly split
evenly between categorical and continuous. Despite the promise of deep
learning to alleviate the burden of engineering features by hand, the nature
of our raw data does not easily lend itself to be input directly into
feedforward neural networks. We still expend considerable engineering
resources transforming user and video data into useful features. The main
challenge is in representing a temporal sequence of user actions and how
these actions relate to the video impression being scored.

We observe that the most important signals are those that describe a user's
previous interaction with the item itself and other similar items, matching
others' experience in ranking ads~\cite{he2014practical}. As an example,
consider the user's past history with the channel that uploaded the video
being scored --- how many videos has the user watched from this channel?
When was the last time the user watched a video on this topic? These
continuous features describing past user actions on related items are
particularly powerful because they generalize well across disparate items.
We have also found it crucial to propagate information from candidate
generation into ranking in the form of features, e.g.\ which sources
nominated this video candidate? What scores did they assign?

Features describing the frequency of past video impressions are also
critical for introducing ``churn'' in recommendations (successive requests
do not return identical lists). If a user was recently recommended a video
but did not watch it then the model will naturally demote this impression on
the next page load. Serving up-to-the-second impression and watch history is
an engineering feat onto itself outside the scope of this paper, but is
vital for producing responsive recommendations.

\textit{Embedding Categorical Features}

Similar to candidate generation, we use embeddings to map sparse categorical
features to dense representations suitable for neural networks. Each unique
ID space (``vocabulary'') has a separate learned embedding with dimension
that increases approximately proportional to the logarithm of the number of
unique values. These vocabularies are simple look-up tables built by passing
over the data once before training. Very large cardinality ID spaces (e.g.\
video IDs or search query terms) are truncated by including only the top $N$
after sorting based on their frequency in clicked impressions.
Out-of-vocabulary values are simply mapped to the zero embedding. As in
candidate generation, multivalent categorical feature embeddings are
averaged before being fed in to the network.

Importantly, categorical features in the same ID space also share underlying
embeddings. For example, there exists a single global embedding of video IDs
that many distinct features use (video ID of the impression, last video ID
watched by the user, video ID that ``seeded'' the recommendation, etc.).
Despite the shared embedding, each feature is fed separately into the
network so that the layers above can learn specialized representations per
feature. Sharing embeddings is important for improving generalization,
speeding up training and reducing memory requirements. The overwhelming
majority of model parameters are in these high-cardinality embedding spaces
--- for example, one million IDs embedded in a 32 dimensional space have 7
times more parameters than fully connected layers 2048 units wide.

\textit{Normalizing Continuous Features}

Neural networks are notoriously sensitive to the scaling and distribution of
their inputs~\cite{ioffe2015batch} whereas alternative approaches such as
ensembles of decision trees are invariant to scaling of individual features.
We found that proper normalization of continuous features was critical for
convergence. A continuous feature $x$ with distribution $f$ is transformed
to $\tilde{x}$ by scaling the values such that the feature is equally
distributed in $[0, 1)$ using the cumulative distribution,
$\tilde{x} = \int_{-\infty}^{x} df$. This integral is approximated with
linear interpolation on the quantiles of the feature values computed in a
single pass over the data before training begins.

In addition to the raw normalized feature $\tilde{x}$, we also input powers
$\tilde{x}^2$ and $\sqrt{\tilde{x}}$, giving the network more expressive
power by allowing it to easily form super- and sub-linear functions of the
feature. Feeding powers of continuous features was found to improve offline
accuracy.

\subsection{Modeling Expected Watch Time}

Our goal is to predict expected watch time given training examples that are
either positive (the video impression was clicked) or negative (the
impression was not clicked). Positive examples are annotated with the amount
of time the user spent watching the video. To predict expected watch time we
use the technique of weighted logistic regression, which was developed for
this purpose.

The model is trained with logistic regression under cross-entropy loss
(Figure~7). However, the positive (clicked) impressions are weighted by the
observed watch time on the video. Negative (unclicked) impressions all
receive unit weight. In this way, the odds learned by the logistic
regression are $\frac{\sum T_i}{N - k}$ where $N$ is the number of training
examples, $k$ is the number of positive impressions, and $T_i$ is the watch
time of the $i$th impression. Assuming the fraction of positive impressions
is small (which is true in our case), the learned odds are approximately
$E[T](1 + P)$, where $P$ is the click probability and $E[T]$ is the expected
watch time of the impression. Since $P$ is small, this product is close to
$E[T]$. For inference we use the exponential function $e^x$ as the final
activation function to produce these odds that closely estimate expected
watch time.

\subsection{Experiments with Hidden Layers}

Table~1 shows the results we obtained on next-day holdout data with
different hidden layer configurations. The value shown for each
configuration (``weighted, per-user loss'') was obtained by considering both
positive (clicked) and negative (unclicked) impressions shown to a user on a
single page. We first score these two impressions with our model. If the
negative impression receives a higher score than the positive impression,
then we consider the positive impression's watch time to be \textit{mispredicted
watch time}. Weighted, per-user loss is then the total amount mispredicted
watch time as a fraction of total watch time over heldout impression pairs.

These results show that increasing the width of hidden layers improves
results, as does increasing their depth. The trade-off, however, is server
CPU time needed for inference. The configuration of a 1024-wide ReLU
followed by a 512-wide ReLU followed by a 256-wide ReLU gave us the best
results while enabling us to stay within our serving CPU budget.

For the $1024 \rightarrow 512 \rightarrow 256$ model we tried only feeding
the normalized continuous features without their powers, which increased
loss by 0.2\%. With the same hidden layer configuration, we also trained a
model where positive and negative examples are weighted equally.
Unsurprisingly, this increased the watch time-weighted loss by a dramatic
4.1\%.

\section{Conclusions}

We have described our deep neural network architecture for recommending
YouTube videos, split into two distinct problems: candidate generation and
ranking.

Our deep collaborative filtering model is able to effectively assimilate
many signals and model their interaction with layers of depth, outperforming
previous matrix factorization approaches used at
YouTube~\cite{weston2011wsabie}. There is more art than science in selecting
the surrogate problem for recommendations and we found classifying a future
watch to perform well on live metrics by capturing asymmetric co-watch
behavior and preventing leakage of future information. Withholding
discrimative signals from the classifier was also essential to achieving
good results --- otherwise the model would overfit the surrogate problem and
not transfer well to the homepage.

We demonstrated that using the age of the training example as an input
feature removes an inherent bias towards the past and allows the model to
represent the time-dependent behavior of popular of videos. This improved
offline holdout precision results and increased the watch time dramatically
on recently uploaded videos in A/B testing.

Ranking is a more classical machine learning problem yet our deep learning
approach outperformed previous linear and tree-based methods for watch time
prediction. Recommendation systems in particular benefit from specialized
features describing past user behavior with items. Deep neural networks
require special representations of categorical and continuous features which
we transform with embeddings and quantile normalization, respectively.
Layers of depth were shown to effectively model non-linear interactions
between hundreds of features.

Logistic regression was modified by weighting training examples with watch
time for positive examples and unity for negative examples, allowing us to
learn odds that closely model expected watch time. This approach performed
much better on watch-time weighted ranking evaluation metrics compared to
predicting click-through rate directly.
